\section{Síntesis y ampliación}

\subsection{Repaso de los puntos clave}

\begin{enumerate}
    \item Los \textbf{modelos de lenguaje por difusión} rompen la restricción de generación izquierda a derecha de los modelos AR, permitiendo generar texto de manera paralela.
    \item El \textbf{MDLM} simplifica el entrenamiento de la difusión discreta haciéndolo tan sencillo como en los modelos AR mediante un ELBO continuo y conciso, logrando un rendimiento a gran escala que puede igualar o incluso superar al de los modelos AR.
    \item La \textbf{memoria KV} es el mayor cuello de botella para la implementación real del MDLM: la atención bidireccional impide congelar los valores de activación.
    \item El \textbf{SoLM}, mediante un diseño que combina ordenamiento y atención causal, mantiene la capacidad de generación paralela mientras soporta la memoria KV, logrando una aceleración en el razonamiento de varios órdenes de magnitud.
    \item En tareas con datos discretos sin sesgo izquierda-derecha (generación de proteínas, moléculas), las ventajas de los modelos de lenguaje por difusión son especialmente notables.
\end{enumerate}

\subsection{Estado actual de la aplicación industrial}

Hasta el momento del grabado de esta conferencia (otoño de 2025), el marco MDLM ha sido ampliamente adoptado en la industria:

\begin{itemize}
    \item \textbf{ByteDance Seed Diffusion}: Construido sobre el marco MDLM.
    \item \textbf{Google Gemini Diffusion}: Utiliza métodos de difusión para la generación de texto.
    \item \textbf{Inception Labs Mercury Coder}: Modelo de difusión orientado a la generación de código.
    \item \textbf{NVIDIA}: Emplea MDLM para la generación de moléculas (ámbito del descubrimiento de fármacos).
\end{itemize}

\subsection{Lecturas adicionales}

\begin{itemize}
    \item Sahoo et al., ``Simple and Effective Masked Diffusion Language Models'' (MDLM), NeurIPS 2024.
    \item Sahoo et al., ``Esoteric Language Models'' (SoLM), 2025.
    \item ``Difusión dualidad'': algoritmo de destilación que genera 1000 palabras en 10 pasos.
    \item LaDa: Large Language Diffusion Models (validación en escala de 8B para MDLM).
    \item BD3LM: Block Diffusion Language Models.
    \item GenMol: Generación de moléculas basada en MDLM.
    \item Página del curso KAIST CS492D: \url{https://mhsung.github.io/kaist-cs492d-fall-2025/}
\end{itemize}

%% --- Fin del contenido --- %%

\end{document}
